这里还有两点需要从上面的推导中单独拎出来。
第一，能够配方算出的，是给定 $\mathbf x_0$ 的
$q(\mathbf x_{t-1}\mid\mathbf x_t,\mathbf x_0)$。
训练时原始样本已知，可以用它计算监督；生成时只知道当前 $\mathbf x_t$，
不能把这个带有干净条件的后验直接当作可执行的逆过程。
网络正是在学习用 $\mathbf x_t,t$ 代替未知干净条件的影响。

第二，固定反向方差 $\nu_t^2I$ 后，两个高斯的 KL 中与均值有关的部分是
\[
\frac1{2\nu_t^2}\|\tilde\mu_t-\mu_\theta(\mathbf x_t,t)\|^2.
\]
把前面得到的均值表达代进去，变为
\[
\frac{\beta_t^2}{2\nu_t^2\alpha_t(1-\bar\alpha_t)}
\|\epsilon-\epsilon_\theta(\mathbf x_t,t)\|^2.
\]
因此，最后常用的无权重噪声 MSE 是一种重新加权。
它便于训练，但不能说时间相关系数在数学上都不影响优化。
如果所有时间步共用同一组参数，不同权重会改变各个噪声级争取模型容量的方式。

再回到采样。均值网络算出 $\mu_\theta$ 后，通常还要执行
$\mathbf x_{t-1}=\mu_\theta(\mathbf x_t,t)+\nu_t z$。
这里重新抽取的 $z$，与训练时用来构造 $\mathbf x_t$ 的噪声标签不是同一个对象。
这就把“前向加噪得到训练对”“回归逆过程参数”“执行逆过程生成”三步接上了。
下面换到 NCSN 的分数视角时，这条关系还会出现。\cite{ddpm,theorybook}
